\documentclass[11pt,spanish]{article}
\renewcommand{\sfdefault}{lmss}
\renewcommand{\familydefault}{\rmdefault}
\usepackage[utf8]{inputenc}
\usepackage{geometry}
\geometry{verbose,tmargin=2cm,bmargin=2cm,lmargin=2cm,rmargin=2cm,columnsep=1cm,marginparwidth=1cm}
\usepackage{fancyhdr}
\pagestyle{fancy}
\setlength{\headheight}{13.6pt}
\usepackage[skip=\bigskipamount]{parskip}
\usepackage{array}
\usepackage{microtype}
\usepackage{hyperref}
\usepackage{pifont}
\usepackage{titling}
\setlength{\droptitle}{-2cm}
\usepackage{xcolor}
\usepackage{enumitem}
\usepackage[most]{tcolorbox}
\usepackage{amssymb}

\definecolor{darkpastelblue}{rgb}{0.47, 0.62, 0.8}
\definecolor{amber}{rgb}{1.0, 0.75, 0.0}
\definecolor{rojo7a3c3d}{rgb}{0.478, 0.235, 0.239}
\definecolor{verde587246}{rgb}{0.345, 0.447, 0.275}
\definecolor{antiquefuchsia}{rgb}{0.57, 0.36, 0.51}
\definecolor{pastelgreen}{rgb}{0.82,0.92,0.78}
\definecolor{indiagreen}{rgb}{0.07,0.53,0.03}

\newcommand{\blank}[1][4cm]{\rule{#1}{0.4pt}}
\newcommand{\cuadro}{\(\square\)}
\newcommand{\cf}[2]{\colorbox{#1}{\textcolor{white}{\strut #2}}}

\newtcolorbox{ejercicio}[1]{
colback=antiquefuchsia!5!white,
colframe=antiquefuchsia!85!black,
fonttitle=\bfseries,
title=#1,
breakable}

\newtcolorbox{observacion}[1]{
colback=amber!5!white,
colframe=amber!85!black,
fonttitle=\bfseries,
title=#1,
breakable}

\newtcolorbox{respuesta}[1]{
colback=pastelgreen!35!white,
colframe=indiagreen!85!black,
fonttitle=\bfseries,
title=#1,
breakable}

\setlength{\parskip}{\bigskipamount}
\setlength{\parindent}{0pt}

\AtBeginDocument{
  \def\labelitemi{\(\star\)}
  \def\labelitemii{\(\ast\)}
  \def\labelitemiii{\(\bullet\)}
  \def\labelitemiv{\(\circ\)}
}

\usepackage{babel}
\addto\shorthandsspanish{\spanishdeactivate{~<>.}}
\deactivatequoting

\begin{document}

\lhead{PDS}
\rhead{Cuadernillo 01}
\title{Programación de Sistemas\\
Unidad 1: Introducción a la Programación de Sistemas y los Compiladores\\ \bigskip
\cf{rojo7a3c3d}{\textbf{Cuadernillo de ejercicios 01}}\\
\bigskip
\Large\textbf{Panorama de la programación de sistemas}
\ifrespuestas\\[0.4em]\large Versión para docente\fi}
\author{Manuel Soto Romero \and Araceli Liliana Reyes Cabello}
\date{Semestre 2026-2 \\ Colegio de Ciencia y Tecnología UACM}

\maketitle

Este cuadernillo acompaña la Nota 01. Su propósito es practicar cuatro ideas: qué distingue a la programación de sistemas, cuáles son sus áreas principales, cómo se diferencian compiladores e intérpretes y por qué un compilador puede integrar el recorrido del curso. Las actividades avanzan de reconocer conceptos a utilizarlos en situaciones breves.

\begin{observacion}{Observación 1. Límite de trabajo}
Responde únicamente con lo estudiado en la Nota 01. No necesitas describir las fases internas de un compilador, escribir reglas de \textsc{IMP} ni explicar técnicas de generación de código.
\end{observacion}

\begin{observacion}{Observación 2. Ruta sugerida}
Los ejercicios \textbf{1, 3, 5 y 8 son esenciales}: permiten comprobar las ideas centrales de la sesión. Los ejercicios \textbf{2, 4, 6, 7 y 9 son de extensión}: profundizan, aplican o recuperan esas ideas. Si el tiempo es limitado, completa primero los esenciales; ningún ejercicio se elimina del cuadernillo.
\end{observacion}

\section*{\cf{verde587246}{1. Conceptos fundamentales}}

\begin{ejercicio}{Ejercicio 1. Completar las ideas centrales \hfill {\small [Esencial]}}
Completa las oraciones con las expresiones del banco. Cada expresión se usa una sola vez.

\begin{center}
\small
\begin{tabular}{|c|c|c|}
\hline
interfaces & aplicación & niveles de abstracción \tabularnewline
\hline
programación de sistemas & arquitectura de la computadora & software de sistemas \tabularnewline
\hline
\end{tabular}
\end{center}

\begin{enumerate}[leftmargin=*]
    \item El \blank[5.4cm] apoya la operación y el uso de una computadora.
    \item La \blank[5.4cm] se ocupa del diseño y la implementación de ese software.
    \item Una \blank[5.4cm] se concentra principalmente en resolver una tarea concreta para una persona u organización.
    \item La \blank[6cm] incluye instrucciones, registros y memoria.
    \item La programación de sistemas relaciona distintos \blank[5.5cm], desde representaciones cercanas a las personas hasta otras que una máquina puede procesar.
    \item Las \blank[5cm] permiten que la salida de una herramienta tenga la forma que espera la siguiente.
\end{enumerate}
\end{ejercicio}

\ifrespuestas
\begin{respuesta}{Respuesta 1}
1. software de sistemas; 2. programación de sistemas; 3. aplicación; 4. arquitectura de la computadora; 5. niveles de abstracción; 6. interfaces.
\end{respuesta}
\fi

\ifrespuestas\else\newpage\fi
\begin{ejercicio}{Ejercicio 2. Aplicación o software de sistemas \hfill {\small [Extensión]}}
Escribe \textbf{A} si el caso se enfoca principalmente en una aplicación y \textbf{S} si se enfoca principalmente en software de sistemas. Después justifica los casos 3 y 6.

\begin{center}
\small
\begin{tabular}{|p{0.07\textwidth}|p{0.78\textwidth}|p{0.07\textwidth}|}
\hline
 & \textbf{Caso} & \textbf{A/S} \tabularnewline
\hline
1 & Programa que calcula el promedio de un grupo y presenta el resultado. & \tabularnewline[0.55cm]
\hline
2 & Herramienta que traduce un programa a otra representación. & \tabularnewline[0.55cm]
\hline
3 & Programa que administra los recursos de la computadora. & \tabularnewline[0.55cm]
\hline
4 & Sistema para registrar préstamos de una biblioteca escolar. & \tabularnewline[0.55cm]
\hline
5 & Herramienta que ayuda a observar la ejecución de otro programa para localizar errores. & \tabularnewline[0.55cm]
\hline
6 & Programa que reúne código objeto y bibliotecas para formar un ejecutable. & \tabularnewline[0.55cm]
\hline
\end{tabular}
\end{center}

\textbf{Justificación del caso 3:} \blank[10cm]

\blank[14cm]

\textbf{Justificación del caso 6:} \blank[10cm]

\blank[14cm]
\end{ejercicio}

\ifrespuestas
\begin{respuesta}{Respuesta 2}
Clasificación esperada: 1-A, 2-S, 3-S, 4-A, 5-S y 6-S.

Justificaciones posibles:
\begin{itemize}[leftmargin=*]
    \item \textbf{Caso 3:} es software de sistemas porque administra recursos y proporciona servicios necesarios para otros programas.
    \item \textbf{Caso 6:} es software de sistemas porque prepara la forma ejecutable de otros programas al reunir sus partes.
\end{itemize}
\end{respuesta}
\fi

\ifrespuestas\newpage\fi
\section*{\cf{verde587246}{2. Áreas y herramientas}}

\begin{ejercicio}{Ejercicio 3. Relacionar responsabilidades \hfill {\small [Esencial]}}
Relaciona cada herramienta con su función general. Escribe la letra correspondiente en la tabla de respuestas.

\begin{center}
\small
\begin{tabular}{|p{0.30\textwidth}|p{0.60\textwidth}|}
\hline
\textbf{Herramienta} & \textbf{Función} \tabularnewline
\hline
1. Ensamblador & A. Administra recursos y proporciona servicios para ejecutar programas. \tabularnewline
\hline
2. Ligador & B. Ayuda a crear y modificar el texto de un programa. \tabularnewline
\hline
3. Cargador & C. Traduce lenguaje ensamblador a código objeto o de máquina. \tabularnewline
\hline
4. Procesador de macros & D. Coloca un programa en memoria y lo prepara para su ejecución. \tabularnewline
\hline
5. Sistema operativo & E. Ayuda a observar la ejecución y localizar errores. \tabularnewline
\hline
6. Editor & F. Reúne código objeto y bibliotecas para formar un ejecutable. \tabularnewline
\hline
7. Depurador & G. Reconoce macros y las expande en secuencias del lenguaje correspondiente. \tabularnewline
\hline
\end{tabular}
\end{center}

\begin{center}
\begin{tabular}{|c|c|c|c|c|c|c|}
\hline
1 & 2 & 3 & 4 & 5 & 6 & 7 \tabularnewline
\hline
\phantom{C} & \phantom{F} & \phantom{D} & \phantom{G} & \phantom{A} & \phantom{B} & \phantom{E} \tabularnewline[0.45cm]
\hline
\end{tabular}
\end{center}
\end{ejercicio}

\ifrespuestas
\begin{respuesta}{Respuesta 3}
1-C, 2-F, 3-D, 4-G, 5-A, 6-B y 7-E.
\end{respuesta}
\fi

\ifrespuestas\else\newpage\fi
\begin{ejercicio}{Ejercicio 4. Del código fuente a la memoria \hfill {\small [Extensión]}}
Completa el recorrido panorámico con las expresiones del banco.

\begin{center}
\small
\begin{tabular}{|c|c|c|c|}
\hline
memoria & ligador & código objeto & compilador \tabularnewline
\hline
ejecutable & cargador & código en lenguaje ensamblador & ensamblador \tabularnewline
\hline
\end{tabular}
\end{center}

\begin{center}
\renewcommand{\arraystretch}{1.55}
\begin{tabular}{|p{0.31\textwidth}|p{0.24\textwidth}|p{0.29\textwidth}|}
\hline
\textbf{Entrada} & \textbf{Herramienta} & \textbf{Resultado} \tabularnewline
\hline
código fuente & \blank[3cm] & \blank[4cm] \tabularnewline
\hline
\raggedright código en lenguaje\\ensamblador & \blank[3cm] & \blank[4cm] \tabularnewline
\hline
código objeto y bibliotecas & \blank[3cm] & \blank[4cm] \tabularnewline
\hline
ejecutable & \blank[3cm] & \blank[4cm] \tabularnewline
\hline
\end{tabular}
\end{center}

Explica con una oración la diferencia entre \textbf{lenguaje ensamblador} y \textbf{ensamblador}:

\blank[14cm]

\blank[14cm]
\end{ejercicio}

\ifrespuestas
\begin{respuesta}{Respuesta 4}
Por filas: compilador / código en lenguaje ensamblador; ensamblador / código objeto; ligador / ejecutable; cargador / memoria.

El \textbf{lenguaje ensamblador} es una representación del programa; el \textbf{ensamblador} es la herramienta que traduce esa representación a código objeto o código de máquina.
\end{respuesta}
\fi

\section*{\cf{verde587246}{3. Compiladores e intérpretes}}

\begin{ejercicio}{Ejercicio 5. Comparar dos procesadores de lenguajes \hfill {\small [Esencial]}}
Completa la tabla. Usa cada expresión del banco una vez.

\begin{center}
\small
\begin{tabular}{|c|c|c|}
\hline
resultados de la ejecución & traduce & programa fuente y datos \tabularnewline
\hline
programa destino & ejecuta operaciones & programa fuente \tabularnewline
\hline
\end{tabular}
\end{center}

\begin{center}
\renewcommand{\arraystretch}{1.45}
\begin{tabular}{|p{0.25\textwidth}|p{0.31\textwidth}|p{0.31\textwidth}|}
\hline
 & \textbf{Compilador} & \textbf{Intérprete} \tabularnewline
\hline
Entrada principal & \blank[3.7cm] & \blank[3.7cm] \tabularnewline
\hline
Acción principal & \blank[3.7cm] & \blank[3.7cm] \tabularnewline
\hline
Resultado principal & \blank[3.7cm] & \blank[3.7cm] \tabularnewline
\hline
\end{tabular}
\end{center}
\end{ejercicio}

\ifrespuestas
\begin{respuesta}{Respuesta 5}
\begin{center}
\renewcommand{\arraystretch}{1.25}
\begin{tabular}{|p{0.25\textwidth}|p{0.31\textwidth}|p{0.31\textwidth}|}
\hline
 & \textbf{Compilador} & \textbf{Intérprete} \tabularnewline
\hline
Entrada principal & programa fuente & programa fuente y datos \tabularnewline
\hline
Acción principal & traduce & ejecuta operaciones \tabularnewline
\hline
Resultado principal & programa destino & resultados de la ejecución \tabularnewline
\hline
\end{tabular}
\end{center}
\end{respuesta}
\fi

\begin{ejercicio}{Ejercicio 6. Identificar el modelo \hfill {\small [Extensión]}}
Escribe \textbf{C} para compilador, \textbf{I} para intérprete o \textbf{M} si la situación indica una combinación de ambos modelos.

\begin{center}
\small
\begin{tabular}{|p{0.76\textwidth}|p{0.12\textwidth}|}
\hline
\textbf{Situación} & \textbf{C/I/M} \tabularnewline
\hline
La herramienta recibe un programa fuente y produce un programa destino que después puede ejecutarse varias veces. & \tabularnewline[0.65cm]
\hline
La herramienta recibe el programa fuente y los datos de entrada, y produce directamente el resultado de esa ejecución. & \tabularnewline[0.65cm]
\hline
El sistema primero traduce el programa a una forma intermedia y luego ejecuta esa forma mediante otra herramienta. & \tabularnewline[0.65cm]
\hline
El resultado principal del procesamiento es otro programa equivalente escrito en un lenguaje distinto. & \tabularnewline[0.65cm]
\hline
No se entrega un programa destino independiente; se realizan las operaciones indicadas en el programa fuente. & \tabularnewline[0.65cm]
\hline
\end{tabular}
\end{center}

Elige una situación y justifica tu respuesta:

\blank[14cm]

\blank[14cm]
\end{ejercicio}

\ifrespuestas
\begin{respuesta}{Respuesta 6}
Orden esperado: C, I, M, C, I.

Ejemplo de justificación para la tercera situación: corresponde a una combinación porque primero ocurre una traducción y después una herramienta ejecuta la representación obtenida.
\end{respuesta}
\fi

\section*{\cf{verde587246}{4. El compilador como caso integrador}}

\begin{ejercicio}{Ejercicio 7. Cuatro ideas para construir una herramienta \hfill {\small [Extensión]}}
Relaciona cada decisión con la idea que representa: \textbf{especificar}, \textbf{transformar}, \textbf{dividir responsabilidades} o \textbf{comprobar}.

\begin{center}
\small
\begin{tabular}{|p{0.68\textwidth}|p{0.22\textwidth}|}
\hline
\textbf{Decisión del equipo} & \textbf{Idea} \tabularnewline
\hline
Definir con claridad las reglas del lenguaje que recibirá la herramienta. & \tabularnewline[0.7cm]
\hline
Cambiar la representación del programa sin perder el comportamiento que debe conservarse. & \tabularnewline[0.7cm]
\hline
Organizar el sistema en componentes con entradas y salidas explícitas. & \tabularnewline[0.7cm]
\hline
Probar programas correctos, revisar errores y verificar el resultado producido. & \tabularnewline[0.7cm]
\hline
\end{tabular}
\end{center}
\end{ejercicio}

\ifrespuestas
\begin{respuesta}{Respuesta 7}
En orden: especificar, transformar, dividir responsabilidades y comprobar.
\end{respuesta}
\fi

\begin{ejercicio}{Ejercicio 8. El caso del compilador del curso \hfill {\small [Esencial]}}
El producto acumulativo del curso puede representarse así:

\begin{center}
\fbox{programa en \textsc{IMP} o \textsc{IMP++}}
\(\longrightarrow\)
\fbox{compilador del curso}
\(\longrightarrow\)
\fbox{programa en \textsc{C}}
\end{center}

Responde brevemente.

\begin{enumerate}[leftmargin=*]
    \item ¿Qué lenguaje base construiremos en notas y clase? \blank[4.5cm]
    \item ¿Qué extensión desarrollaremos en las prácticas? \blank[4.5cm]
    \item ¿Cuál es el lenguaje destino? \blank[5cm]
    \item ¿Por qué este compilador es software de sistemas?

    \ifrespuestas\else
    \blank[14cm]

    \blank[14cm]
    \fi

    \item Un equipo genera un programa en \textsc{C}, pero nunca revisa si conserva el comportamiento esperado. ¿Cuál de las cuatro ideas del ejercicio anterior está descuidando y por qué?

    \ifrespuestas\else
    \blank[14cm]

    \blank[14cm]
    \fi

    \item Explica en dos o tres oraciones por qué la construcción de este compilador puede integrar distintos aprendizajes del curso.

    \ifrespuestas\else
    \blank[14cm]

    \blank[14cm]

    \blank[14cm]
    \fi
\end{enumerate}
\end{ejercicio}

\ifrespuestas
\begin{respuesta}{Respuesta 8}
\begin{enumerate}[leftmargin=*]
    \item \textsc{IMP}.
    \item \textsc{IMP++}.
    \item \textsc{C}.
    \item Porque procesa otros programas: recibe un programa fuente, detecta errores y produce un programa destino.
    \item \textbf{Comprobar}, porque debe verificarse que el programa generado conserve el comportamiento esperado.
    \item Integra la especificación de un lenguaje, la transformación entre representaciones, la organización de componentes y la comprobación de resultados y errores.
\end{enumerate}
\end{respuesta}
\fi

\newpage
\section*{\cf{verde587246}{5. Revisión de afirmaciones}}

\begin{ejercicio}{Ejercicio 9. Detectar y corregir confusiones \hfill {\small [Extensión]}}
Cada afirmación contiene un error. Subraya la parte incorrecta y reescribe la oración de forma adecuada.

\begin{enumerate}[leftmargin=*]
    \item Un compilador y un intérprete hacen exactamente lo mismo porque ambos producen siempre un programa destino independiente.

    \blank[14cm]

    \blank[14cm]

    \item El cargador reúne código objeto y bibliotecas para formar un ejecutable.

    \blank[14cm]

    \blank[14cm]

    \item El lenguaje ensamblador es la herramienta que traduce código a lenguaje de máquina.

    \blank[14cm]

    \blank[14cm]

    \item La programación de sistemas se limita al desarrollo de sistemas operativos.

    \blank[14cm]

    \blank[14cm]
\end{enumerate}
\end{ejercicio}

\ifrespuestas
\begin{respuesta}{Respuesta 9}
Correcciones esperadas:
\begin{enumerate}[leftmargin=*]
    \item Un compilador traduce un programa fuente y produce un programa destino; un intérprete ejecuta las operaciones indicadas sin producir como resultado principal un programa destino independiente.
    \item El \textbf{ligador} reúne código objeto y bibliotecas para formar un ejecutable; el cargador coloca el ejecutable en memoria.
    \item El \textbf{lenguaje ensamblador} es una representación del programa; el \textbf{ensamblador} es la herramienta que la traduce a código objeto o de máquina.
    \item La programación de sistemas incluye sistemas operativos, ensambladores, ligadores, cargadores, procesadores de macros, compiladores, intérpretes, editores y depuradores, entre otras herramientas.
\end{enumerate}
\end{respuesta}
\fi

\end{document}
